% --- ARCHIVO: Capitulos/composicion_party.tex ---

\chapter{Composición del Ejército}

\begin{loreblock}
``Sesenta puntos. Ese es tu mundo. Construye sabiamente, comandante — con ellos ganarás o perderás antes de que comience la batalla.''
\end{loreblock}

\vspace{0.4cm}

\section{Formato Estándar: 60 Puntos}

El único formato actualmente disponible es el \textbf{Estándar}. Cada jugador tiene \textbf{60 puntos} para construir su ejército (llamado \textit{lista}).

\begin{reglaespecial}{Reglas de Construcción de Lista}
\begin{itemize}
    \item Solo pueden incluirse unidades de \textbf{una única facción}.
    \item Debes consultar el \textbf{Codex de tu facción} para equipar a tus héroes con Equipo, Artefactos y Magias.
    \item \textbf{No puedes superar los 60 puntos}, pero puedes presentar una lista de menos de 60 como válida.
    \item Al armar tu lista, asigna el equipamiento antes de la partida. Los ítems no pueden cambiarse una vez que comienza.
\end{itemize}
\end{reglaespecial}

\section{Espacios de Equipamiento por Unidad}

Cada unidad tiene una serie de ranuras fijas de equipamiento. No todas las unidades pueden ocupar todas las ranuras — esto varía según la raza, clase y características propias de la unidad (indicadas en su ficha del Codex).

\begin{tabularx}{\linewidth}{|l|X|}
\hline
\rowcolor{HammerRed!15} \textbf{Zona del Cuerpo} & \textbf{Ranuras Disponibles} \\ \hline
\textbf{Cabeza} & Casco, Hombreras \\ \hline
\textbf{Brazos} & Guantes, Mangas, Anillo 1, Anillo 2 \\ \hline
\textbf{Piernas} & Pantalón, Calzado \\ \hline
\textbf{Cuerpo} & Pechera, Cinturón, Collar \\ \hline
\textbf{Armas} & 2×1 Mano / 1×2 Manos / 1 Mano + Adicional (artefacto u otra arma) \\ \hline
\textbf{Spells} & 1 a 6 ranuras (según la unidad) \\ \hline
\textbf{Gema} & Puede o no equipar una Gema (según la unidad) \\ \hline
\textbf{Inventario} & 1 a 6 ranuras (según la unidad) \\ \hline
\end{tabularx}

\vspace{0.3cm}

El \textbf{inventario} puede contener consumibles (pociones, pergaminos, ítems de uso único). Cada ítem consumible indica cuándo puede ser usado y qué acción requiere.

\begin{imagebox}{Diagrama visual de las ranuras de equipamiento de una unidad — silueta con flechas señalando cada zona de equipo}
    \vspace{3.5cm}
\end{imagebox}

\section{Tipos de Armadura}

La armadura determina la Reducción de Daño (RD) base de la unidad. Existen tres categorías:

\begin{reglaespecial}{Categorías de Armadura}
\begin{itemize}
    \item \textbf{Túnica (Ligera):} No otorga RD, pero proporciona bonificadores mágicos. Ideal para hechiceros.
    \item \textbf{Armadura Mediana:} Proporciona RD moderada sin penalizar el movimiento. Ideal para guerreros ágiles y DPS.
    \item \textbf{Armadura Pesada:} Máxima RD, pero \textbf{reduce el movimiento} del portador según la penalización indicada en la pieza. Ideal para tanques y unidades de primera línea.
\end{itemize}
\end{reglaespecial}

\section{El Héroe: Unidad Líder}

Cada ejército puede incluir un \textbf{Héroe} o \textbf{Líder}, que es la unidad más poderosa de la lista. El Líder:

\begin{itemize}
    \item Es el único que puede portar la \textbf{Gema} de la facción.
    \item Tiene acceso a equipamiento exclusivo de alto nivel.
    \item Tiene un costo en puntos significativamente mayor al de las unidades regulares.
    \item Genera el \textbf{Área de Influencia} para distribución de PE.
\end{itemize}

\begin{tcolorbox}[colback=GoldRule!10, colframe=GoldRule, title=\textbf{REGLA: Elegido / Líder}]
Si el Líder es eliminado, el Área de Influencia desaparece y ninguna unidad puede usar la Gema por el resto de la partida (a menos que otra unidad pueda robarla y activarla según las reglas de su facción).
\end{tcolorbox}

\section{La Ficha de Unidad}

Cada unidad en juego tiene una ficha que resume sus estadísticas, equipamiento y habilidades. La ficha se divide en:

\begin{itemize}
    \item \textbf{Frente:} Estadísticas base, localización de vida, imagen de la unidad y equipamiento actual.
    \item \textbf{Dorso:} Descripción completa de poderes, habilidades especiales y texto de lore.
\end{itemize}

\begin{imagebox}{Ejemplo de ficha de unidad completa (frente y dorso) — formato imprimible, tamaño carta o A4}
    \vspace{4cm}
\end{imagebox}

\subsection*{Ficha Genérica de Unidad (Resumen)}

\begin{tabularx}{\linewidth}{|l|c|X|}
\hline
\rowcolor{HammerRed!15} \textbf{Atributo} & \textbf{Valor} & \textbf{Descripción} \\ \hline
\textbf{P. Acción} & 1/2/3/4 & Acciones posibles por ronda \\ \hline
\textbf{P. Energía} & 1/2/3/4 & Casteos / habilidades mágicas por ronda \\ \hline
\textbf{Robar} & +\_\_ & Bono a tirada de 1d12 para interactuar/robar \\ \hline
\textbf{Focus} & +\_\_ & Bono para extender efectos/magia \\ \hline
\textbf{Movimiento} & \_\_ cm & Distancia por cada 1 PA gastado \\ \hline
\textbf{Melee ATK} & +\_\_ & Bono a tirada de 1d12 en Melee \\ \hline
\textbf{Range ATK} & +\_\_ & Bono a tirada de 1d12 a Distancia \\ \hline
\textbf{Reflejos} & +\_\_ & Dificultad para ser impactado (esquiva) \\ \hline
\textbf{Bloqueo} & +\_\_ & Dificultad para ser impactado (defensa) \\ \hline
\textbf{Res. Física} & +\_\_ & Bono a 1d6 contra estados físicos \\ \hline
\textbf{Res. Mágica} & +\_\_ & Bono a 1d6 contra estados mágicos \\ \hline
\end{tabularx}

\vspace{0.4cm}

\begin{reglaespecial}{Regla de Vínculo de Gema}
En la ficha aparece el campo: \textbf{Vínculo Gema: [SÍ] [NO]}. Este marcador debe estar visible durante la partida e indica si la unidad tiene acceso activo a la energía de la Gema de su Líder.
\end{reglaespecial}
